\paragraph{Mode and provenance of the conditional branches.}
The saved Stage-I, II, and III configurations specify \texttt{robust\_incremental}. In the recovered Stage-III dispatcher this calls the robust adaptive-box core and reuses previously scanned momentum intervals; it is not a different thermodynamic objective. The Stage-IV launch script requests the same mode with basin clustering and selective refinement, at most three refined minima, three optional basins, and a rank-and-cap overflow policy. These are launch-script settings, not proof of every historical shell override. The export does not bind every invocation to an executable hash. The following branch descriptions identify the recovered implementation; unrecorded call-specific settings remain unconfirmed.

For the saved two-dimensional configurations, the condensation-energy ambiguity tolerance is $\epsilon_F=10^{-6}t$, the strictly-positive-amplitude comparison tolerance is $\epsilon_+=10^{-8}t$, and $\Delta_\epsilon=10^{-3}t$. They serve different purposes. In the robust core, candidate basins are ordered by their refined objective value, not by phase preference; the near-degeneracy window is $\max(10^{-5}t,5\epsilon_F)$. The dispatcher requests at least 800 momentum and 600 amplitude samples for local refinement. The historical maximum expansion count, absolute momentum bound, and optional basin-selection switches are not fixed by these run JSON files alone and are not inferred from current defaults.

\paragraph{Momentum-window diagnostics and expansion.}
The recovered core examines the first and last three momentum cells, reduced for very short grids. Edge sensitivity requires a positive-amplitude trace or an objective below $-\epsilon_F$. It flags an edge within $5\times10^{-5}t$ of the scanned minimum, a low-energy local minimum near that edge, an outward descent, or an optimum within the recorded edge margin. The left and right flags determine directional expansion; an unbracketed case can request one symmetric expansion. On each affected side the padding is $\max(4|\delta q|,W_q/2)$, clipped at the supplied absolute momentum bound; the new grid preserves or improves the old spacing. No new interval is added when clipping prevents growth. Thus a diagnostic enclosure test, a termination bound, and the returned unresolved flag are distinct quantities.

\paragraph{Near-zero amplitude and returned states.}
The guard is triggered when enabled and either the best objective is within $\epsilon_F$ of zero, the amplitude is within the supplied refinement half-width of $\Delta_\epsilon$, or it lies below $\Delta_\epsilon$. A separate scan over amplitudes starting at $\Delta_\epsilon$ provides the stored positive-amplitude comparison. For a normal candidate, a comparison below $-\epsilon_+$ with amplitude above $\Delta_\epsilon$ replaces the candidate; a value in $[0,\epsilon_+]$ marks a normal-side boundary band. Robust superconducting acceptance requires both amplitude above $\Delta_\epsilon$ and objective below $-\epsilon_F$; otherwise the recovered branch assigns normal or ambiguous status according to the recorded comparison.

In this core, the returned momentum-unresolved bit specifically requires a superconducting candidate, an optimum at the final edge margin, and exhaustion of the permitted expansions; it is not identical to every edge-risk diagnostic. Nonfinite selected momentum, amplitude, or objective sets a failure bit. A zero status code with no ambiguity is trusted; a normal-side boundary-band record can remain training-eligible without being trusted. Rerun-required is the negation of training eligibility. These implementation distinctions prevent a numerical status from silently becoming a fourth phase. Full per-call binding of optional refinement switches remains a reproducibility limitation.

\paragraph{Cutoff and zero-temperature record.}
The recovered shared integration kernel uses $|t\cos k-\mu|<\omega_D$ before the momentum sum and retains the configured full-grid denominator. With the stated cutoff and sampled parameter ranges this mask does not remove momentum points. The exported occupation helper evaluates the sigmoid expression directly and contains no explicit zero-temperature branch. A separately version-bound zero-temperature guard was not recovered; no limiting prescription is supplied here as though it were a recorded historical setting. This is an implementation-provenance limitation, not a change to or assessment of the objective above.

\paragraph{Online spectral status versus offline qualification.}
The recovered online attachment first requires a superconducting thermodynamic candidate, which is broader than the FFLO-only subset displayed in the main figures. It scans $[-\pi,\pi]$ including both endpoints and compares the minimum absolute eigenvalue with $\max(\epsilon_g^{\rm abs},\epsilon_g^{\rm rel}E_*)$. Nonfinite spectral data or an unresolved Pfaffian sign returns unresolved status; a gap at or below the threshold is recorded as gapless, not as a gapped topological phase. Spectral trust additionally requires training eligibility and no momentum/amplitude unresolved bits. The Stage-IV launch script supplies default grid size 2048, relative gap tolerance $10^{-8}$ and absolute tolerance zero, with environment overrides possible. These online defaults do not certify an offline refined gap. The separately stored topology-pass implementation doubles selected gap grids and locally refines small-gap/small-margin cases; that script targets a different dataset and its settings are not automatically attributed to the Fig.3 endpoint. Its invocation-level applicability remains to be confirmed.
